\documentclass[11pt]{article}

\usepackage[T1]{fontenc}
\usepackage[utf8]{inputenc}
\usepackage[a4paper,margin=1in]{geometry}
\usepackage{amsmath,amssymb,amsthm,mathtools}
\usepackage[hidelinks]{hyperref}

\newtheorem{theorem}{Theorem}
\newtheorem{lemma}{Lemma}

\newcommand{\Z}{\mathbb{Z}}
\newcommand{\tomega}{\vec{\omega}}
\newcommand{\domn}{\operatorname{dom}}
\newcommand{\dic}{\operatorname{dic}}
\newcommand{\AC}{\operatorname{AC}}
\newcommand{\Cay}{\operatorname{Cay}}

\title{The \(k=3\) Case of Conjecture 5.10}
\author{}
\date{}

\begin{document}

\maketitle

\begin{theorem}
For every odd integer \(n=2m+1\geq 7\), the almost-consecutive
circulant tournament
\[
  \AC_n
  =
  \Cay\!\left(
    \Z/n\Z,\,
    \{1,\ldots,m-1\}\cup\{m+1\}
  \right),
\]
where
\[
  i\longrightarrow j
  \quad\Longleftrightarrow\quad
  j-i\pmod n\in g,
  \qquad
  g=\{1,\ldots,m-1\}\cup\{m+1\},
\]
is \(3\)-\(\tomega\)-critical.

Consequently, there are infinitely many \(3\)-\(\tomega\)-critical
tournaments, so the \(k=3\) case of Conjecture~5.10 of
Aboulker--Aubian--Charbit--Lopes~\cite{AACL} holds. Moreover,
Question~5.9 of \cite{AACL} has a negative answer for \(k=3\).
\end{theorem}

Recall that the tournament clique number is
\[
  \tomega(T)
  =
  \min_{\sigma}
  \omega\!\left(B_{\sigma}(T)\right),
\]
where the minimum is over all total orderings \(\sigma\) of \(V(T)\),
and \(B_{\sigma}(T)\) is the back-edge graph: if
\(i<_{\sigma}j\), then \(\{i,j\}\) is an edge of
\(B_{\sigma}(T)\) exactly when \(j\to i\) is an arc of \(T\).
A tournament \(T\) is \(k\)-\(\tomega\)-critical if
\[
  \tomega(T)=k
  \quad\text{and}\quad
  \tomega(T-v)=k-1
  \quad\text{for every }v\in V(T).
\]

\begin{proof}
We divide the proof into four steps.

\medskip
\noindent
\textbf{Step 1: \(\AC_n\) is a tournament.}
For each \(r\in\{1,\ldots,m\}\), exactly one of \(r\) and \(n-r\)
belongs to \(g\). Indeed, if \(r\leq m-1\), then \(r\in g\), whereas
\[
  n-r=2m+1-r\geq m+2
\]
does not belong to \(g\). For \(r=m\), we have \(m\notin g\), while
\(n-m=m+1\in g\). Thus \(g\) contains exactly one element from each
pair \(\{r,-r\}\) of nonzero residues, so \(\AC_n\) is a tournament.
It is vertex-transitive, since every rotation \(x\mapsto x+c\) is an
automorphism. Notice that primality of \(n\) is not used.

\medskip
\noindent
\textbf{Step 2: the upper bound \(\tomega(\AC_n)\leq 3\).}
Consider the identity order
\[
  0<1<\cdots<n-1.
\]
For \(0\leq i<j\leq n-1\), the pair \(\{i,j\}\) is a back edge if and
only if \(j\to i\), or equivalently
\[
  i-j\pmod n\in g.
\]
In terms of the ordinary integer gap \(j-i\), this is equivalent to
\[
  j-i\in
  D:=\{m\}\cup\{m+2,\ldots,2m\}.
\]
In particular, every back-edge gap is at least \(m\).

Suppose that
\[
  a_1<a_2<\cdots<a_q
\]
is a clique in the back-edge graph. Each consecutive gap is at least
\(m\), and therefore
\[
  a_q-a_1\geq(q-1)m.
\]
On the other hand, \(a_q-a_1\leq n-1=2m\). Hence \(q\leq 3\), proving
\[
  \tomega(\AC_n)\leq 3.
\]

We also record the equality case. If
\(a_1<a_2<a_3\) is a back-edge triangle, then both consecutive gaps
are at least \(m\), while \(a_3-a_1\leq2m\). Consequently both gaps
equal \(m\), so
\[
  \{a_1,a_2,a_3\}
  =
  \{a_1,a_1+m,a_1+2m\}.
\]
Since \(a_1+2m\leq2m\), we must have \(a_1=0\). Thus the unique
back-edge triangle in the identity order is
\[
  \{0,m,2m\}.
\]

\medskip
\noindent
\textbf{Step 3: the lower bound \(\tomega(\AC_n)\geq3\).}
Let
\[
  N_0=N^+[0]
  =
  \{0\}\cup g
  =
  \{0,1,\ldots,m-1,m+1\}.
\]
Then \(|N_0|=m+1\), and vertex-transitivity gives
\[
  N^+[v]=N_0+v
  \qquad\text{for every }v\in\Z/n\Z.
\]
A pair \(\{0,v\}\) dominates \(\AC_n\) if and only if
\[
  N_0\cup(N_0+v)=\Z/n\Z.
\]
Since \(2|N_0|-n=1\), this is equivalent to
\[
  |N_0\cap(N_0+v)|=1.
\]
No singleton dominates, because \(|N_0|=m+1<n\). The required
domination lower bound therefore follows from the next lemma.

\begin{lemma}\label{lem:overlap}
For every odd integer \(n=2m+1\geq7\),
\[
  \min_{t\neq0}|N_0\cap(N_0+t)|=2.
\]
\end{lemma}

\begin{proof}
Put
\[
  J=\{0,1,\ldots,m+1\},
  \qquad
  N_0=J\setminus\{m\},
\]
and let
\[
  K=(\Z/n\Z)\setminus J=\{m+2,\ldots,2m\}.
\]
Thus \(|J|=m+2\) and \(|K|=m-1\). By taking complements,
\[
  |J\cap(J+t)|
  =
  n-|K\cup(K+t)|
  =
  3+|K\cap(K+t)|.
\]
For \(1\leq t\leq m\), the set \(K\) is an interval of length
\(m-1\), and hence
\[
  |K\cap(K+t)|=\max\{0,m-1-t\}.
\]
It follows that
\[
  |J\cap(J+t)|
  =
  \begin{cases}
    m+2-t, & 1\leq t\leq m-1,\\
    3,     & t=m.
  \end{cases}
\]

Passing from \(J\) to \(N_0=J\setminus\{m\}\) removes \(m\) from the
intersection whenever \(m-t\in J\), which holds for every
\(1\leq t\leq m\). It also removes \(m+t\) whenever \(m+t\in J\),
which occurs exactly when \(t=1\). Therefore
\[
  |N_0\cap(N_0+t)|
  =
  \begin{cases}
    m-1,   & t=1,\\
    m+1-t, & 2\leq t\leq m-1,\\
    2,     & t=m.
  \end{cases}
\]
All these quantities are at least \(2\), since \(m\geq3\).
Finally,
\[
  |N_0\cap(N_0+t)|
  =
  |N_0\cap(N_0-t)|,
\]
so the calculation for \(1\leq t\leq m\) covers every nonzero
residue \(t\). The minimum is therefore \(2\).
\end{proof}

For completeness, when \(n=7\), every nonzero \(t\) attains the
minimum in Lemma~\ref{lem:overlap}. For every odd \(n\geq9\), the
minimizers are exactly
\[
  \{m-1,m,m+1,m+2\}.
\]

Lemma~\ref{lem:overlap} implies that no pair dominates \(\AC_n\), and
hence
\[
  \domn(\AC_n)\geq3.
\]
Property~3.2 of \cite{AACL} states that
\[
  \domn(T)\leq\tomega(T)\leq\dic(T)
\]
for every tournament \(T\). Consequently,
\[
  \tomega(\AC_n)\geq3.
\]
Together with Step~2, this proves
\[
  \tomega(\AC_n)=3.
\]

\medskip
\noindent
\textbf{Step 4: criticality.}
By vertex-transitivity, it suffices to delete the vertex \(0\). The
unique back-edge triangle in the identity order is \(\{0,m,2m\}\).
Therefore the restriction of the identity order to
\(V(\AC_n)\setminus\{0\}\) has a triangle-free back-edge graph, and
so
\[
  \tomega(\AC_n-0)\leq2.
\]

Conversely, \(\AC_n-0\) contains the directed triangle
\[
  1\longrightarrow2\longrightarrow m+3\longrightarrow1.
\]
Indeed, the corresponding differences are
\[
  1,\qquad m+1,\qquad
  1-(m+3)\equiv m-1\pmod n,
\]
all of which belong to \(g\). The vertices \(1,2,m+3\) are distinct
and nonzero, since \(m\geq3\). A tournament has clique number \(1\)
if and only if it is transitive, so the presence of this directed
triangle gives
\[
  \tomega(\AC_n-0)\geq2.
\]
Thus
\[
  \tomega(\AC_n-v)=2
  \qquad\text{for every }v\in V(\AC_n).
\]
This proves that \(\AC_n\) is \(3\)-\(\tomega\)-critical.

Distinct odd values of \(n\) give tournaments of distinct orders.
Hence
\[
  \{\AC_n:n\geq7\text{ odd}\}
\]
is an infinite family of pairwise nonisomorphic
\(3\)-\(\tomega\)-critical tournaments.

It remains to state explicitly why Question~5.9 fails for \(k=3\).
Every proper subtournament \(A\subsetneq\AC_n\) omits some vertex
\(v\), and therefore
\[
  A\subseteq\AC_n-v.
\]
By monotonicity under taking subtournaments,
\[
  \tomega(A)
  \leq
  \tomega(\AC_n-v)
  =
  2.
\]
Thus the only subtournament of \(\AC_n\) with clique number at least
\(3\) is \(\AC_n\) itself. Since \(n\) is unbounded, no bound
\(\ell(3)\) as in Question~5.9 can exist.
\end{proof}

\section*{Relation to the result of Crew--Fan--Koerts--Moore--Spirkl}

There is no contradiction with the 2026 result of
Crew--Fan--Koerts--Moore--Spirkl~\cite{CFKMS}. Their theorem confirms
Conjecture~5.8 of \cite{AACL}: there exist functions \(f\) and
\(\ell\) such that
\[
  \tomega(T)\geq f(k)
  \quad\Longrightarrow\quad
  \text{\(T\) contains \(X\) with
  \(|V(X)|\leq\ell(k)\) and \(\tomega(X)\geq k\)}.
\]
Question~5.9 asks for the strictly stronger statement in which the
threshold function is the identity, namely \(f(k)=k\). The
tournaments \(\AC_n\) all satisfy \(\tomega(\AC_n)=3\), so they
disprove this identity-threshold statement at \(k=3\). They do not
conflict with a theorem whose hypothesis at \(k=3\) may require
\(\tomega(T)\geq f(3)>3\).

\section*{Provenance and dependencies}

The construction, the identity-order upper bound, the uniqueness of
the back-edge triangle, Lemma~\ref{lem:overlap}, and the criticality
argument are self-contained. The sole external input is
Property~3.2 of \cite{AACL}, used to pass from
\(\domn(\AC_n)\geq3\) to \(\tomega(\AC_n)\geq3\).

\begin{thebibliography}{2}

\bibitem{AACL}
P.~Aboulker, G.~Aubian, P.~Charbit, and R.~Lopes,
\newblock \emph{Clique number of tournaments},
\newblock arXiv:2310.04265, 2023.

\bibitem{CFKMS}
L.~Crew, X.~Fan, H.~Koerts, B.~Moore, and S.~Spirkl,
\newblock \emph{Characterizing large clique number in tournaments},
\newblock arXiv:2602.09863, 2026.

\end{thebibliography}

\end{document}
